\chapter{Hidden Operators}
\label{chap:Hidden_Operators}

\section{Purpose}

Hidden Operators are latent transformations SEGO may apply to symbolic
structures that are not obvious from surface syntax.

Examples:
\begin{itemize}
    \item archetypal substitutions,
    \item dualities (light/shadow, inner/outer),
    \item symmetry operations (swap roles, invert polarity),
    \item structural inversions (inside/outside, above/below).
\end{itemize}

They provide depth and nuance in symbolic reasoning, but must be tightly
controlled.

\section{Operator Model}

An operator $H$ acts on symbolic objects (Engrams, manifolds, traces):
\[
H: X \rightarrow X'
\]

With metadata:
\[
H = \langle id,\, domain,\, effect,\, constraints \rangle
\]

Where:
\begin{itemize}
    \item domain: what it may act on (tags, manifolds, narratives),
    \item effect: transformation description,
    \item constraints: when and how it may be used.
\end{itemize}

\section{Examples}

\subsection{Shadow Projection Operator}

\begin{itemize}
    \item domain: psychological/narrative structures (PSY, THEO, PHIL),
    \item effect: highlight or simulate shadow/aspect inversion,
    \item constraint: never treat output as literal diagnosis; for narrative
          or symbolic exploration only.
\end{itemize}

\subsection{Role Reversal Operator}

\begin{itemize}
    \item domain: SOCI/PSY/ONTO,
    \item effect: swap perspectives between two roles,
    \item constraint: mark as hypothetical; may not be used to rewrite
          C-class history.
\end{itemize}

\section{Governance}

\begin{itemize}
    \item Hidden Operators must be documented in Esoterica Index.
    \item Certain operators may be tied to SIGIL pathways for invocation.
    \item Use in high-impact contexts requires EGO supervision and possibly HITL.
\end{itemize}

\section{Summary}

Hidden Operators are the secret moves of SEGO’s symbolic game. By codifying
and governing them, we keep their power available without letting it become
unbounded or opaque.
